\chapter[What Geometrogenesis Must Derive]{What Geometrogenesis Must Derive}
\label{ch:geometrogenesis}

\paragraph{Learning outcomes.}
After this chapter the reader should be able to:
\begin{enumerate}
  \item state the geometrogenesis problem as a derivation task,
        not a metaphor;
  \item explain why analog-gravity metrics are not general relativity;
  \item identify the representation content that a microscopic order
        parameter must carry if massless spin-2 modes are to exist;
  \item list the consistency conditions treated as non-negotiable in
        this book: diffeomorphism invariance,
        \(\nabla^\mu T_{\mu\nu}^{\mathrm{total}}=0\), a bounded
        infrared Hamiltonian, and an explicit map from microscopic
        couplings to \(\Geff\);
  \item see why Chapter~\ref{ch:preliminaries} is required before any
        later proof can be stated.
\end{enumerate}

\section{What this book is trying to derive}
\label{sec:ch1-goal}

General relativity is an infrared theory of a dynamical Lorentzian
metric \(g_{\mu\nu}\) whose vacuum excitations are two massless,
transverse-traceless spin-2 polarizations, and whose coupling to matter
is fixed by local energy-momentum conservation and the equivalence
principle. In standard presentations the metric is primitive. One writes
an Einstein--Hilbert action
\begin{equation}
  \label{eq:EH}
  S_{\mathrm{EH}}[g]
  =\frac{1}{16\pi G}\int\mathrm{d}^4x\,\sqrt{-g}\,R
\end{equation}
and proceeds.

Geometrogenesis asks a more primitive question:

\begin{quote}
Can \(g_{\mu\nu}\), the Einstein--Hilbert dynamics, and the two graviton
polarizations be obtained as the long-wavelength collective description
of a microscopic quantum system that does not contain a metric as a
fundamental field?
\end{quote}

That question is the subject of this book. The framework developed here
is Non-Abelian Matrix-Condensate Geometrogenesis (NAMCG). The
microscopic object is not a single complex scalar. It is a
matrix-valued order parameter \(\Phi^I{}_A(x)\) carrying an internal
gauge index \(I\) and an emergent tangent-space index
\(A\in\{0,1,2,3\}\). The metric is a composite,
\begin{equation}
  \label{eq:composite-metric}
  e^A_\mu(x)\equiv\langle\partial_\mu\Phi^{IA}(x)\rangle,\qquad
  g_{\mu\nu}(x)=\eta_{AB}\,e^A_\mu(x)\,e^B_\nu(x).
\end{equation}

The later chapters argue that entanglement equilibrium on causal
diamonds forces the infrared mass of the symmetric trace-free shear
mode to vanish, that a heat-kernel expansion supplies an explicit
formula for \(\Geff\), and that the entanglement backreaction stress
tensor remains covariantly conserved.

Those are theorems \emph{inside the framework}. This first chapter does
not prove them. It makes precise what would count as a proof, and why
the obvious scalar shortcuts fail.

\section{Analog gravity, and the ceiling it cannot break}
\label{sec:ch1-analog}

A large and useful literature constructs effective metrics for
perturbations on condensed-matter backgrounds. The standard example is
a weakly interacting Bose gas with order parameter
\(\Psi=\sqrt{\rho}\,e^{i\theta}\). Phase fluctuations \(\theta\) obey,
after linearization, a wave equation
\begin{equation}
  \label{eq:acoustic-wave}
  \frac{1}{\sqrt{-g_{\mathrm{ac}}}}
  \partial_\mu\Bigl(
    \sqrt{-g_{\mathrm{ac}}}\,
    g_{\mathrm{ac}}^{\mu\nu}\partial_\nu\theta
  \Bigr)=0
\end{equation}
in an acoustic metric \(g_{\mathrm{ac}}^{\mu\nu}\) built from the
background density and flow. Horizons, ergoregions, and Hawking-like
spectra can be engineered in the laboratory.

This is analog gravity. It is not geometrogenesis of general
relativity, for four structural reasons.

\paragraph{Only one scalar mode.}
The Goldstone field of a single complex scalar is a scalar. The
acoustic metric is a rewriting of that scalar wave equation. It does
not create two independent tensor polarizations.

\paragraph{Screening.}
Gross--Pitaevskii interactions introduce a healing length \(\xi\). The
static potential felt by a localized source is of Yukawa or
modified-Bessel type, not a \(1/r\) Newtonian potential with radiative
\(1/r\) tails.

\paragraph{Wrong infrared gravity.}
A relativistic scalar condensate coupled to an auxiliary metric
generically yields Nordstr\"om gravity (conformal, with no bending of
light in the pure-scalar sector) or scalar--tensor gravity. The Einstein
tensor is not selected as the unique linearized kinetic operator for a
massless spin-2 field.

\paragraph{Dissipation inserted by hand.}
Binary inspiral in general relativity follows from quadrupole radiation
in a theory with massless spin-2 waves. Replacing that derivation by a
Peters formula or a drag term in Newton's equation may fit a curve. It
does not derive the curve from the microscopic action.

NAMCG is written to climb over that ceiling, not to decorate it.

\begin{figure}[ht]
  \centering
  \begin{tikzpicture}[x=1.05cm,y=1.05cm,>=Stealth,font=\small]
  % ---- left panel: acoustic analog ----
  \begin{scope}[shift={(0,0)}]
    \draw[thick] (0,0) rectangle (5.2,4.6);
    \node[font=\bfseries] at (2.6,4.25) {scalar condensate};
    % background flow
    \foreach \y in {0.55,1.15,1.75,2.35} {
      \draw[-{Stealth[length=1.6mm]},gray] (0.35,\y) -- (2.15,\y);
    }
    \node[gray,font=\scriptsize] at (1.25,2.75) {flow $\mathbf{v}$};
    % source
    \fill (1.25,1.45) circle (0.07);
    % acoustic cones (tilted by flow)
    \draw[thick] (1.25,1.45) -- (2.35,2.25);
    \draw[thick] (1.25,1.45) -- (2.35,0.65);
    \draw[thick] (1.25,1.45) -- (2.55,1.45);
    \node[font=\scriptsize,align=center] at (3.7,2.55) {acoustic\\cone};
    \node[font=\scriptsize,align=left] at (3.55,1.15) {one scalar\\mode $\theta$};
    \node[font=\scriptsize] at (2.6,0.28) {analog metric $g_{\mathrm{ac}}$};
  \end{scope}

  % ---- arrow ----
  \draw[-{Stealth[length=2.4mm]},very thick] (5.45,2.3) -- (6.15,2.3);

  % ---- right panel: GR / NAMCG target ----
  \begin{scope}[shift={(6.4,0)}]
    \draw[thick] (0,0) rectangle (5.2,4.6);
    \node[font=\bfseries] at (2.6,4.25) {emergent Lorentzian $g_{\mu\nu}$};
    % light cone
    \draw[thick] (2.6,0.7) -- (1.15,3.55);
    \draw[thick] (2.6,0.7) -- (4.05,3.55);
    \draw[thick] (2.6,0.7) -- (2.6,3.7);
    \fill (2.6,0.7) circle (0.06);
    \node[font=\scriptsize] at (2.6,0.42) {$p$};
    % two TT marks
    \node[font=\scriptsize,align=center] at (1.05,3.95) {$+$};
    \node[font=\scriptsize,align=center] at (4.15,3.95) {$\times$};
    \draw[thick] (0.7,3.15) -- (1.4,3.15);
    \draw[thick] (1.05,2.8) -- (1.05,3.5);
    \draw[thick] (3.75,3.35) -- (4.35,2.75);
    \draw[thick] (3.75,2.75) -- (4.35,3.35);
    \node[font=\scriptsize,align=center] at (2.6,2.15)
      {two TT modes\\helicities $\pm 2$};
    \node[font=\scriptsize] at (2.6,0.18) {massless spin-2};
  \end{scope}
\end{tikzpicture}
  \caption{A scalar order parameter supplies an acoustic metric.
    It does not supply the two massless spin-2 polarizations of
    general relativity.}
  \label{fig:scalar-vs-gr}
\end{figure}

\section{Representation theory in one paragraph}
\label{sec:ch1-reps}

Local Lorentz invariance acts on tangent-space indices. A graviton is a
symmetric, traceless, divergenceless tensor \(h_{\mu\nu}\) with two
on-shell degrees of freedom in four dimensions. A single scalar
\(\phi\) can produce a conformal mode proportional to
\(g_{\mu\nu}\delta\phi\), and perhaps a vector current if a global
symmetry is gauged, but it cannot produce an irreducible shear tensor
as an independent propagating excitation of the order parameter itself.

That is the representation-theoretic reason the microscopic field is
promoted to \(\Phi^I{}_A\). The tangent index \(A\) is not cosmetic. It
is the handle on which local \(\mathrm{SO}(3,1)\) will later act.
Chapter~5 decomposes
\begin{equation}
  \label{eq:lorentz-split}
  \delta\Phi^I{}_A
  =\frac14\eta_{\mu\nu}\delta g^{\mu\nu}\,I
  +A_\mu\gamma^\mu
  +T_{\mu\nu}
\end{equation}
into conformal, vector, and symmetric trace-free pieces. Chapter~7 is
the argument that \(T_{\mu\nu}\) becomes massless.

If that argument fails, NAMCG collapses back to analog gravity. The
book is organized so the reader can see exactly where it would fail.

\section{What must be delivered before one may say ``this is GR''}
\label{sec:ch1-checklist}

A derivation of infrared general relativity from a microscopic model is
not complete unless all of the following are obtained, not assumed.

\begin{enumerate}
  \item \textbf{Composite metric.}
    \(g_{\mu\nu}\) is defined from microscopic expectation values. It
    is not put into the fundamental action as an independent Einstein
    metric.
  \item \textbf{Local Lorentz frames.}
    The vacuum must lock internal indices to tangent indices so that
    \(e^A_\mu\) transforms as a vierbein and no physical Goldstone
    scalars remain. That is Chapter~10.
  \item \textbf{Two massless TT modes.}
    The infrared kinetic operator of \(T_{\mu\nu}\) must be the
    Lichnerowicz / Einstein operator, with \(m_T^2=0\). That is
    Chapters~6--7.
  \item \textbf{Coefficient map.}
    Matching the physical \(G_{\mathrm{eff}}\) requires specified
    gauge-fixed operators, determinant weights, sector thresholds,
    and a scheme. Chapter~8 gives a conditional truncated identity;
    denote the continued curvature coefficient for those inputs by
    \(C_R^L\); its physical value remains open in Problem~F:
    \begin{equation}
      \label{eq:geff-preview}
      C_R^L[\{\mathcal O_s,w_s,{\LUV}_s,\mu_s\},\mathcal S]
      \quad\text{is not yet evaluated.}
    \end{equation}
  \item \textbf{Bianchi consistency.}
    Whatever extra stress is attributed to entanglement or condensate
    backreaction must satisfy
    \begin{equation}
      \label{eq:conservation-preview}
      \nabla^\mu\bigl(
        T^{\mathrm{matter}}_{\mu\nu}
        +T^{\mathrm{tensor}}_{\mu\nu}
        +\Tent_{\mu\nu}
      \bigr)=0.
    \end{equation}
    That is Chapter~11. Without it, one has a non-conservative force,
    not geometry.
  \item \textbf{Ultraviolet control.}
    Higher-derivative terms generated by integrating out the condensate
    cannot leave propagating Ostrogradsky ghosts in the infrared. That
    is Chapter~9.
  \item \textbf{Vacuum energy.}
    The same heat-kernel expansion produces a term
    \(c_0\LUV^4\). A framework that claims to recover Einstein's
    equation must say why this is not a Planck-scale cosmological
    constant. That is Chapter~12.
  \item \textbf{Phenomenology as corollary.}
    Rotation curves, regular cores, inspiral drag, and tidal disruption
    are allowed only after items 1--7. They are not evidence that
    items 1--7 have been proved.
\end{enumerate}

This checklist is the spine of Parts~II--IV. Parts~V--VI apply it.

\section{The NAMCG derivation in one view}
\label{sec:ch1-spine}

The logical order used from Chapter~3 onward is
\begin{equation}
  \label{eq:spine}
  \begin{aligned}
    \Phi^I{}_A
    &\;\longrightarrow\;
    \langle\Phi\rangle=v\,\delta^I{}_A
    \;\longrightarrow\;
    e^A_\mu=\langle\partial_\mu\Phi^{IA}\rangle \\
    &\;\longrightarrow\;
    g_{\mu\nu}=\eta_{AB}e^A_\mu e^B_\nu
    \;\longrightarrow\;
    \delta\Phi=\text{conformal}+\text{vector}+T_{\mu\nu} \\
    &\;\longrightarrow\;
    \delta S_D=0
    \;\longrightarrow\;
    m_T\to 0
    \;\longrightarrow\;
    S_{\mathrm{EH}}[g]+S_{\mathrm{NAMCG}}^{\mathrm{(backreaction)}}.
  \end{aligned}
\end{equation}

\begin{figure}[ht]
  \centering
  \begin{tikzpicture}[
  x=1cm,y=1cm,
  box/.style={
    draw, thick, rounded corners=1.2pt,
    align=center, font=\scriptsize,
    inner sep=3.4pt, minimum height=0.78cm
  },
  arr/.style={-{Stealth[length=1.7mm]}, thick},
  font=\small
]
  % row 1
  \node[box] (phi) at (0,0)
    {$\Phi^I{}_A$};
  \node[box] (vev) at (2.55,0)
    {$\langle\Phi\rangle=v\delta^I{}_A$};
  \node[box] (e) at (5.35,0)
    {$e^A_\mu=\langle\partial_\mu\Phi^{IA}\rangle$};
  \draw[arr] (phi) -- (vev);
  \draw[arr] (vev) -- (e);

  % row 2
  \node[box] (g) at (0,-1.45)
    {$g_{\mu\nu}=\eta_{AB}e^A_\mu e^B_\nu$};
  \node[box] (split) at (4.15,-1.45)
    {$\delta\Phi=\psi+A+T_{\mu\nu}$};
  \draw[arr] (e.south) -- ++(0,-0.28) -| (g.north);
  \draw[arr] (g) -- (split);

  % row 3
  \node[box] (ee) at (0,-2.9)
    {$\delta S_D=0$};
  \node[box] (mt) at (2.7,-2.9)
    {$m_T\to 0$};
  \node[box] (eh) at (5.55,-2.9)
    {$S_{\mathrm{EH}}[g]+S_{\mathrm{br}}$};
  \draw[arr] (split.south) -- ++(0,-0.28) -| (ee.north);
  \draw[arr] (ee) -- (mt);
  \draw[arr] (mt) -- (eh);

  \node[font=\scriptsize,text=black!70] at (2.7,-3.55)
    {Eq.~\eqref{eq:spine}; $S_{\mathrm{br}}=$ backreaction};
\end{tikzpicture}
  \caption{The NAMCG derivation spine used in this book. A single
    scalar phase cannot lock internal indices to local Lorentz frames.}
  \label{fig:spine}
\end{figure}

Two comments are essential already.

First, the entanglement step is doing real work. In a generic
spontaneously broken matrix model, \(T_{\mu\nu}\) may be massive. The
central physical claim of the book is that Jacobson's stationarity
condition on dyadic causal diamonds,
\begin{equation}
  \label{eq:deltaS-preview}
  \delta S_D=\delta\langle H_D\rangle=0,
\end{equation}
is not an extra thermodynamic slogan. It is a Ward identity of the
microscopic theory, and that identity forbids \(m_T^2\neq 0\) in the
infrared.

Second, the entropic drag must not be mistaken for a derived result.
The inspiral integrator of Chapter~15 is a reduced Newtonian
laboratory for the phenomenological law
\begin{equation}
  \label{eq:drag-preview}
  \mathbf{F}_{\mathrm{diss}}=-\eta\,\frac{\mu}{d^4}\,\mathbf{v}.
\end{equation}
Its bookkeeping does not establish that this force is the reduction
of a covariantly conserved NAMCG stress. That microscopic connection
remains open in Problem~\ref{sec:status-I}. If the reader treats the
Python drag as the theory, the book has failed.

\section{How the later phenomenology sits on the spine}
\label{sec:ch1-pheno}

Once the infrared Einstein sector exists, NAMCG adds three controlled
deformations, each with a named chapter.

\paragraph{Holographic Boundary-Entropy Vacuum (HBEV).}
A boundary functional involving \(K/8\pi G\) and an entropy density
\(S_{\mathrm{ent}}(h_{ij},\phi)\) is used to regulate
\(\Lambda_{\mathrm{eff}}\). This is how the book addresses the
\(c_0\LUV^4\) term rather than ignoring it.

\paragraph{Einsteinian cubic gravity.}
A ghost-controlled cubic curvature invariant motivates a
regular-core model for compact objects (a phenomenological profile,
not an established exact solution),
\begin{equation}
  \label{eq:fr-preview}
  f(r)=1-\frac{2Mr^2}{r^3+r_0^3}+\frac{\Lambda_{\mathrm{eff}}}{3}r^2,
\end{equation}
so that \(f(r)\to 1+\Lambda_{\mathrm{eff}}r^2/3\) as \(r\to 0\).

\paragraph{Vacuum backreaction in galaxies and compact binaries.}
The same conserved extra tensor sources a modification of Poisson's
equation; a phenomenological drag in compact binaries is tested as a
separate assumption, not derived from that tensor.

None of these is a substitute for Chapters~3--11. They are tests of
whether the extra tensors have the right long-distance falloff and the
right conservation law.

\section{How to read this book}
\label{sec:ch1-howto}

Chapters~2--5 are definitional. Read them with a pen. If a symbol is
used later with two meanings, the fault is in those chapters and should
be caught there.

Chapters~6--11 are the proof core. Each theorem is written with
hypotheses first. If a hypothesis is an assumption of the framework
rather than a lemma, it is labeled as such.

Chapters~12--15 are applications. Numerical sections flag reduced
models in the first paragraph.

Chapter~16 restates the whole derivation as a single chain. Chapter~17
lists what would falsify the framework. Read Chapter~17 before deciding
that the theory is closed.

\paragraph{Prerequisites.}
\paragraph{What you need.}
Special relativity, the Euler--Lagrange equation, and the Einstein
equation as a field equation for a metric. Linear algebra at the
level of commutators and traces. Quantum field theory is helpful
but is introduced at the point of use. Entanglement entropy is
built in Chapter~6 and Appendix~B. If a later proof needs more
than that, the extra tool is installed in Chapter~2 or in an
appendix, not presupposed as coursework.

\section{Exercises}
\label{sec:ch1-exercises}

\begin{enumerate}
  \item \textbf{Acoustic metric.}
    Starting from the continuity and Bernoulli equations of an
    irrotational barotropic fluid, derive \(g_{\mathrm{ac}}^{\mu\nu}\)
    and show that linearized phonons see that metric. Where does a
    healing length appear if the fluid is a Gross--Pitaevskii
    condensate?
  \item \textbf{Mode count.}
    In \(3+1\) dimensions, how many on-shell degrees of freedom does a
    massive symmetric tensor have? A massless one with gauge invariance
    \(\delta h_{\mu\nu}=\partial_\mu\xi_\nu+\partial_\nu\xi_\mu\)?
    Why does that counting matter for Section~\ref{sec:ch1-reps}?
  \item \textbf{Scalar-tensor ceiling.}
    Consider
    \(S=\int\sqrt{-g}\bigl(\phi R-(\partial\phi)^2-V(\phi)\bigr)\).
    Linearize about Minkowski spacetime and a constant \(\phi_0\).
    Show that the propagating excitations are a massless spin-2 field
    mixed with a scalar. What extra condition would be needed to freeze
    the scalar?
  \item \textbf{Conservation as a constraint.}
    Suppose a two-body equation
    \(\mu\ddot{\mathbf{r}}=-\nabla U+\mathbf{a}_{\mathrm{irrev}}\) is
    written in Minkowski space. What property must
    \(\mathbf{a}_{\mathrm{irrev}}\) have if there is to exist a stress
    tensor \(\Tent_{\mu\nu}\) with
    \(\partial^\mu(T^{\mathrm{matter}}_{\mu\nu}+\Tent_{\mu\nu})=0\)?
    A full curved-space answer is Chapter~11; a Newtonian answer is
    enough here.
  \item \textbf{Checklist audit.}
    Take any analog-gravity paper you know and score it against the
    eight items of Section~\ref{sec:ch1-checklist}. This exercise is
    the standard the book applies to itself.
\end{enumerate}

\section{Bridge to Chapter~2}
\label{sec:ch1-bridge}

Every later proof is a sentence in a language that has not yet been
fixed: manifolds and vierbeine, Einstein--Hilbert variation and Bianchi
identities, heat kernels, causal diamonds, and modular Hamiltonians.
Chapter~1 stated the derivation problem. Chapter~2 installs the
language at the exact thickness later chapters consume---no more,
because a reader who must first finish a complete treatise on general
relativity will never reach the matrix condensate.

The first object Chapter~2 defines is the one this chapter kept using
without ceremony: a Lorentzian manifold with a non-degenerate metric of
signature \((-,+,+,+)\), and the difference between a primitive metric
and a composite vierbein bilinear.